%% 郭宸毓 — 履歷（繁體中文）  [v4：日期改為單行區間；內容重新配重，不再一面倒寫地震]
%% 編譯：tectonic -X compile cv-zh.tex
%%
%% 版面使用 Weitian LI 的 resume.cls（LPPL 1.3c），已原封不動置於本目錄：
%%   https://github.com/liweitianux/resume
%% 其本身衍生自 Christophe Roger 的 YACC 與 Plasmati Graduate CV。

\documentclass[zh]{resume}

%% 類別預設繁簡不分（Noto Serif CJK SC），改用繁體字形；
%% 同時稍微縮小字級與行距，讓內容落在兩頁。
\setmainfont{IBM Plex Serif}[Scale=0.93]
\setmonofont{IBM Plex Mono}[Scale=0.87]
\setCJKmainfont{Noto Serif CJK TC}[Scale=0.93]

\usepackage{geometry}
\geometry{top=1.0cm, bottom=0.9cm, left=1.05cm, right=1.05cm, includefoot}
\linespread{1.06}
\setlength{\parskip}{0.35ex plus 0.3ex minus 0.1ex}
\setlist[itemize,1]{label=\faAngleRight, leftmargin=1.2em, labelsep=0.45em,
                    topsep=0.2ex, itemsep=0.25ex, parsep=0pt}
\titlespacing{\section}{0em}{0.6ex}{0.6ex}

%% v4：日期改寫成單行區間（「2023/11 – 2024/11」），不再用類別預設的上下堆疊寫法
%% ——堆疊會被誤讀成兩個獨立日期。最寬的區間量起來是 86.3pt，日期欄要從 60pt 加到 90pt；
%% 右欄退還同樣的 30pt，總寬不變。
\addtolength{\leftcolwidth}{30pt}
\addtolength{\rightcolwidth}{-30pt}

\fileinfo{%
  \faGithub{} \link{https://github.com/whes1015/cv}{whes1015/cv} \hspace{0.5em}
  版面：\link{https://github.com/liweitianux/resume}{liweitianux/resume}（LPPL~1.3c）\hspace{0.5em}
  \faEdit{} \today
}

\name{宸毓}{郭}
\keywords{Go, TypeScript, Python, C++, Rust, Flutter, eBPF, XDP, CoreDNS, 嵌入式, ESP32,
  PyTorch, LLM 量化, 地震學, 地震預警}
\tagline{Chen-Yu Kuo \enspace\textbullet\enspace 系統與全端工程師 --- eBPF／Go 基礎設施、ESP32 韌體、
  機器學習、LLM 推論與地震預警}

\profile{
  \email{whes1015@gmail.com}
  \github{whes1015}
  \iconlink{\faRobot}{https://huggingface.co/YuYu1015}{\texttt{huggingface.co/YuYu1015}} \\
  \degree{電機工程學系 學士}
  \university{國立高雄科技大學（NKUST）}
  \address{高雄，臺灣}
}

\begin{document}
\makeheader

%======================================================================
\begin{abstract}
我習慣把系統從頭做到尾，而且多半是唯一的作者：以 Go 寫的 eBPF/XDP 封包過濾器與橫跨五個城市的
Nginx／CoreDNS 節點群、逐暫存器親手撰寫的 ESP32 韌體、Flutter 與 TypeScript 用戶端，
以及自己訓練、評估並量化的 PyTorch 模型。多數成果掛在我創辦、並擔任兩位管理者之一的
\textbf{ExpTech（探索科技）}之下——29 人的志工工程組織，維運臺灣規模最大的獨立地震預警平臺。
此外曾任衛波科技後端工程師、中央研究院資訊科學研究所兼任助理，現於國家衛生研究院從事程式開發。
\end{abstract}

%======================================================================
\sectionTitle{技術能力}{\faWrench}
%======================================================================
\begin{competences}[6em]
  \comptence{程式語言}{%
    TypeScript／JavaScript、Go、Python、Dart、C/C++、Rust、Lua、Shell
  }
  \comptence{基礎設施}{%
    Linux、eBPF/XDP、Docker、Nginx/OpenResty、CoreDNS、NTP/SNTP、Prometheus、CI/CD
  }
  \comptence{後端與應用}{%
    Go 與 Node／TypeScript 服務、REST／WebSocket／SSE、Flutter（iOS／Android）、Electron、MapLibre
  }
  \comptence{嵌入式}{%
    ESP32/Arduino、FreeRTOS、I\textsuperscript{2}C/SPI、LSM6DS3、GNSS（QZSS L1S）、LoRa/Meshtastic
  }
  \comptence{機器學習}{%
    PyTorch、SeisBench、ONNX、scikit-learn、vLLM、llm-compressor、TensorRT Model Optimizer
  }
  \comptence{領域知識}{%
    地震預警、GMPE 校正、中央氣象署／JMA 震度階、地震波相位辨識
  }
  \comptence{\icon{\faLanguage} 語言}{%
    \textbf{中文}——母語；\textbf{英文}——技術文件讀寫
  }
\end{competences}

%======================================================================
\sectionTitle{學歷}{\faGraduationCap}
%======================================================================
\begin{educations}
  %% TODO：把「現在」改成預計畢業時間，例如「2022 – 2027/06」
  \education%
    {}[2022 -- 現在]%
    {國立高雄科技大學}%
    {電機工程學系}%
    {電機工程}%
    {學士}
\end{educations}

%======================================================================
\sectionTitle{工作經歷}{\faBriefcase}
%======================================================================
\begin{experiences}
  %% 整段起訖寫在「結束日期」欄位、不傳起始日期：resume.cls 只有在後面接了描述或標籤時
  %% 才會收掉表格列，單寫 \experience[起]{迄}{摘要} 會和下一列黏在一起。
  \experience%
    {2025/01 -- 現在}%
    {\textbf{程式開發} --- 國家衛生研究院（NHRI），臺灣}

  \separator{0.6em}

  \experience%
    {2024/12 -- 2025/06}%
    {\textbf{兼任助理} --- 中央研究院 資訊科學研究所，臺北}

  \separator{0.6em}

  \experience%
    {2023/11 -- 2024/11}%
    {\textbf{後端工程師} --- 衛波科技股份有限公司（P-Waver Inc.，地震預警系統），臺北}

  \separator{0.6em}

  \experience%
    {2020/10 -- 現在}%
    {\textbf{創辦人暨技術負責人} --- ExpTech（探索科技），臺灣}%
    [\begin{itemize}
      \item 擔任組織兩位管理者之一，帶領 29 人團隊、234 個儲存庫，涵蓋韌體、37 個 Go 服務、
        Flutter 用戶端與機器學習流程。於組織內提交 \textbf{13,411} 次 commit
        （GitHub 全站 19,284 次、477 個 PR），並在所有核心專案皆為最大貢獻者——
        DPIP 838/2,339、TREM-Lite 448/847、平臺 API 480/803——另獨力完成 665 次 commit 的
        正式環境基礎設施儲存庫。
      \item 2022 年底與\textbf{中央氣象署}簽訂資料合作——《親子天下》報導其為地震測報中心
        第一組非公司組織的合作對象；並以合作者身分進出測報中心逾 10 次（《天下雜誌》，2026 年 2 月）。
      \item 面向真實使用者、承擔真實風險：DPIP 於 Google Play 累計 \textbf{10 萬次以上}安裝、
        App Store 1,378 則評分平均 4.2；TREM 桌面版安裝檔累計 \textbf{328,542} 次下載。
        2024 年 4 月 3 日花蓮 M7.2 地震當日，觀測網於一日內記錄到 800 起以上地震事件。
    \end{itemize}]
\end{experiences}

%======================================================================
\sectionTitle{資安、開源與媒體}{\faLock}
%======================================================================
\begin{itemize}
  \item \textbf{資安與事件應變}——架設蜜罐，於 \textbf{CVE-2025-55182}
    （React Server Components，CVSS 10.0）2025 年 12 月揭露後數日內捕獲實際在野的攻擊行為；
    記錄並驗證 6 個惡意程式樣本雜湊值與 15 個攻擊來源 IP。
  \item \textbf{開源貢獻}——於自身組織之外累計 12 個已合併的 PR：為開源地震儀專案 \textbf{AnyShake}
    （\link{https://github.com/anyshake/spectrogram-js}{\texttt{spectrogram-js}}）提交繪製效能優化、
    為 \texttt{smc.peering.tw} 將 Leaflet 遷移至 MapLibre，以及為母校官方學生應用程式
    \textbf{NKUST-AP-Flutter} 升級 Flutter 工具鏈。
  \item \textbf{媒體報導}——《天下雜誌》2026 年 2 月（於中央氣象署地震測報中心人物專訪中具名受訪）；
    《親子天下》2024 年 7 月（DPIP 團隊專題，0403 地震後累計下載約 50 萬次）。
\end{itemize}

\newpage
%======================================================================
\sectionTitle{精選技術成果}{\faCode}
%======================================================================
\begin{projects}
  \project%
    {}[2025/12 -- 現在]%
    {獨力開發}%
    {\,Go、eBPF/XDP、OpenResty、CoreDNS}%
    {正式環境基礎設施——封包過濾、邊緣節點與 DNS}%
    {\begin{itemize}
      \item \textbf{kekkai}——以 Go 搭配 \texttt{cilium/ebpf} 實作封包過濾：XDP 入向程式與 TC 出向掛鉤，
        以 \textsc{lru\_hash}/\textsc{array}/\textsc{ringbuf} map 實現單一 IP 速率限制、
        允許／封鎖名單與連接埠政策，並具緊急卸載路徑。
      \item \textbf{Nginx/OpenResty 伺服器群}——跨臺北、高雄、臺南、東京、大阪節點共 23 個 upstream，
        並以 Lua 模組處理流量統計、NTP 與影像代理。
      \item \textbf{分流式（split-horizon）CoreDNS}，自製 Go \texttt{dnsstats} 外掛（\texttt{miekg/dns}），
        提供各網域／用戶端查詢統計與 Prometheus 指標。另獨力開發平臺的 NTP 伺服器、API 閘道
        與 Meshtastic LoRa 資料接收服務。
    \end{itemize}}

  \separator{0.7em}

  \project%
    {}[2026/04 -- 2026/07]%
    {獨力開發}%
    {\,Python、PyTorch、vLLM、TensorRT Model Optimizer}%
    {大型語言模型推論與模型壓縮}%
    {\begin{itemize}
      \item 於 Hugging Face 發布 \textbf{27 個量化模型}（累計約 \textbf{73,000} 次下載），
        供 vLLM 與 llama.cpp 部署：以 TensorRT Model Optimizer 產出 \textbf{NVFP4}，
        另有 INT4-AutoRound 與 GGUF，涵蓋至 35B 規模的 MoE 架構。
      \item 實作\textbf{拒答方向消融（refusal-direction ablation）}——由成對提示求出拒答方向，
        並自權重矩陣正交化移除，而非在推論時做引導——發布 \emph{Ornith-1.0} 9B/35B 系列，
        含 DPO 微調版本。
    \end{itemize}}

  \separator{0.7em}

  \project%
    {}[2025/09 -- 2026/06]%
    {獨力開發，7,759 行}%
    {\,C++、ESP32、FreeRTOS}%
    {ES-Net —— 感測器韌體與 200 個節點的裝置網路}%
    {\begin{itemize}
      \item 不依賴原廠函式庫，\textbf{逐暫存器手寫 LSM6DS3 加速度計驅動}——自建暫存器對應表，
        涵蓋 \textsc{who\_am\_i}、\textsc{ctrl1\_xl}、\textsc{int1\_ctrl} 與 FIFO 路徑；
        另於裝置端以 IIR 濾波器組搭配樹狀陣列（BIT）即時維護振幅百分位。
      \item 分散式感測裝置最難的是時間：SNTP 採用 \textsc{smooth} 同步模式，讓時鐘平滑校正而非在
        波形中途跳動；並實作軟硬體看門狗、OTA 更新，以及具備 TOTP 裝置驗證的自訂二進位通訊協定。
      \item 目前運行於全臺 \textbf{200 個以上節點}；第一代裝置為自行組裝，單台成本低於新臺幣 800 元。
    \end{itemize}}

  \separator{0.7em}

  \project%
    {}[2023/08 -- 現在]%
    {最大貢獻者，838/2,339 commit}%
    {\,Dart/Flutter、TypeScript}%
    {DPIP —— 跨平臺行動應用程式與後端}%
    {\begin{itemize}
      \item iOS 與 Android 用戶端，實作即時推播告警、MapLibre 圖層與 10 種介面語言；
        Google Play 累計 \textbf{10 萬次以上}安裝、App Store 1,378 則評分平均 4.2。
        獨力撰寫該儲存庫的測試套件與自訂 commit 檢查工具鏈，並完整開發 TypeScript 後端服務
        \texttt{dpip-server-ts}（122/122 commit）。
    \end{itemize}}

  \separator{0.7em}

  \project%
    {}[2026/08 -- 現在]%
    {獨力開發，約 12,000 行、含單元測試}%
    {\,C++}%
    {QZSS L1S 衛星訊息解碼器}%
    {\begin{itemize}
      \item 完整解碼 QZSS L1S 250 位元 SBAS 訊框——前導碼、6 位元訊息型別、212 位元酬載與 CRC-24Q——
        包含型別 43 \textbf{DC Report}（日本氣象廳災防廣播，IS-QZSS-DCR）與型別 44 DCX。
    \end{itemize}}

  \separator{0.7em}

  \project%
    {}[2026/03 -- 2026/07]%
    {獨力開發}%
    {\,Python、PyTorch、SeisBench、ONNX}%
    {應用機器學習 —— 感測資料的迴歸與序列模型}%
    {\begin{itemize}
      \item 以中央氣象署 2000--2025 年共 \textbf{45,398} 筆 PGA/PGV/震度觀測重新校正迴歸模型，
        取代對規模過度放大的舊式公式，並修正場址反應項重複計算（原高估約 2 倍）。
        留出事件驗證：\textbf{完全命中 62.0\%（原生產環境模型 39.7\%）}、
        MAE \textbf{0.400（原 0.671）}、誤報率 0\%。
      \item 訓練 50\,Hz PhaseNet 序列模型，建立 6,457 行的完整流程（資料匯入、雜訊注入、評估、
        ONNX 匯出）。並診斷修復一項視窗設定錯誤：來源視窗短於模型視窗，導致模型學成
        「事件恆在視窗中央」，推論時每 45\,秒 stride 即產生一次誤判。
    \end{itemize}}

  \separator{0.7em}

  \project%
    {}[2023/01 -- 現在]%
    {最大貢獻者，448/847 commit}%
    {\,JavaScript、Electron、Rust}%
    {TREM —— 桌面平臺與第三方外掛生態}%
    {\begin{itemize}
      \item 跨平臺 Electron 用戶端，安裝檔累計 \textbf{328,542} 次下載，跨 74 個發行版本自 2023 年
        1 月起持續發布。設計供第三方擴充的外掛系統並維護其註冊表（417 次 commit）；
        另以 Rust 開發健康監控服務 \texttt{trem-monitor}（105/150）。
    \end{itemize}}
\end{projects}

\end{document}
